% from ega1-10-fr.tex
\oldpage[I]{204}

\subsection{Faisceaux cohérents sur les préschémas formels.}
\label{subsection:I.10.11-fr}

\begin{proposition}[10.11.1]
\label{I.10.11.1-fr}
Si $\mathfrak{X}$ est un préschéma formel localement noethérien, le
faisceau d'anneaux $\mathscr{O}_{\mathfrak{X}}$ est cohérent et tout
faisceau d'idéaux de définition de $\mathfrak{X}$ est cohérent.
\end{proposition}

La question étant locale, on est ramené au cas d'un schéma formel affine
noethérien, et la proposition résulte donc de
(\hyperref[I.10.10.3-fr]{10.10.3}) et
(\hyperref[I.10.10.5-fr]{10.10.5}).

\begin{env}[10.11.2]
\label{I.10.11.2-fr}
Soient $\mathfrak{X}$ un préschéma formel localement noethérien,
$\mathscr{J}$ un faisceau d'idéaux de définition de $\mathfrak{X}$,
$X_n$ le préschéma (usuel) localement noethérien
$(\mathfrak{X},\mathscr{O}_{\mathfrak{X}}/\mathscr{J}^{n+1})$, de
sorte que $\mathfrak{X}$ est limite inductive de la suite $(X_n)$ pour
les morphismes canoniques $u_{mn}:X_m\to X_n$
(\hyperref[I.10.6.3-fr]{10.6.3}). Avec ces notations~:
\end{env}

\begin{theorem}[10.11.3]
\label{I.10.11.3-fr}
Pour qu'un $\mathscr{O}_{\mathfrak{X}}$-Module $\mathscr{F}$ soit
cohérent, il faut et il suffit qu'il soit isomorphe à une limite
projective d'une suite $(\mathscr{F}_n)$, où $\mathscr{F}_n$ est un
$\mathscr{O}_{X_n}$-Module cohérent, tel que
$u_{mn}^*(\mathscr{F}_n)=\mathscr{F}_m$ pour $m\leq n$
(\hyperref[I.10.6.6-fr]{10.6.6}). Le système projectif
$(\mathscr{F}_n)$ est alors isomorphe au système des
$u_n^*(\mathscr{F})=\mathscr{F}
\otimes_{\mathscr{O}_{\mathfrak{X}}}\mathscr{O}_{X_n}$, $u_n$ étant
le morphisme canonique $X_n\to\mathfrak{X}$.
\end{theorem}

La question étant locale, on est ramené au cas où $\mathfrak{X}$ est
un schéma formel affine noethérien, et le théorème est alors conséquence
de (\hyperref[I.10.10.5-fr]{10.10.5}) et
(\hyperref[I.10.10.6-fr]{10.10.6}).

On peut donc dire que \emph{la donnée d'un
$\mathscr{O}_{\mathfrak{X}}$-Module cohérent équivaut à celle d'un
système projectif $(\mathscr{F}_n)$ de
$\mathscr{O}_{X_n}$-Modules cohérents tels que
$u_{mn}^*(\mathscr{F}_n)=\mathscr{F}_m$ pour $m\leq n$}.

\begin{corollary}[10.11.4]
\label{I.10.11.4-fr}
Si $\mathscr{F}$ et $\mathscr{G}$ sont deux
$\mathscr{O}_{\mathfrak{X}}$-Modules cohérents, on peut (avec les
notations de (\hyperref[I.10.11.3-fr]{10.11.3})) définir un
isomorphisme canonique fonctoriel
\[
  \label{I.10.11.4.1-fr}
  \operatorname{Hom}_{\mathscr{O}_{\mathfrak{X}}}
  (\mathscr{F},\mathscr{G})
  \xrightarrow{\sim}
  \varprojlim_n
  \operatorname{Hom}_{\mathscr{O}_{X_n}}
  (\mathscr{F}_n,\mathscr{G}_n).
  \tag{10.11.4.1}
\]
\end{corollary}

La limite projective du second membre doit s'entendre pour les
applications
$\theta_n\mapsto u_{mn}^*(\theta_n)$ ($m\leq n$) de
$\operatorname{Hom}_{\mathscr{O}_{X_n}}
(\mathscr{F}_n,\mathscr{G}_n)$ dans
$\operatorname{Hom}_{\mathscr{O}_{X_m}}
(\mathscr{F}_m,\mathscr{G}_m)$. L'homomorphisme
(\hyperref[I.10.11.4.1-fr]{10.11.4.1}) fait correspondre à un élément
$\theta\in\operatorname{Hom}_{\mathscr{O}_{\mathfrak{X}}}
(\mathscr{F},\mathscr{G})$ la suite $(u_n^*(\theta))$~; on voit
aussitôt qu'on en définit un homomorphisme réciproque du précédent en
faisant correspondre au système projectif
$(\theta_n)\in\varprojlim_n
\operatorname{Hom}_{\mathscr{O}_{X_n}}
(\mathscr{F}_n,\mathscr{G}_n)$
sa limite projective dans
$\operatorname{Hom}_{\mathscr{O}_{\mathfrak{X}}}
(\mathscr{F},\mathscr{G})$, compte tenu de
(\hyperref[I.10.11.3-fr]{10.11.3}).

\begin{corollary}[10.11.5]
\label{I.10.11.5-fr}
Pour qu'un homomorphisme $\theta:\mathscr{F}\to\mathscr{G}$ soit
surjectif, il faut et il suffit que l'homomorphisme correspondant
$\theta_0=u_0^*(\theta):\mathscr{F}_0\to\mathscr{G}_0$ le soit.
\end{corollary}

La question étant locale, on est ramené au cas où
$\mathfrak{X}=\operatorname{Spf}(A)$, $A$ étant adique noethérien,
$\mathscr{F}=M^\Delta$, $\mathscr{G}=N^\Delta$ et $\theta=u^\Delta$,
où $M$ et $N$ sont des $A$-modules de type fini et $u$ un
homomorphisme $M\to N$~; on a en outre alors
$\theta_0=\widetilde{u_0}$, où $u_0$ est l'homomorphisme
$u\otimes1:M\otimes_A A/\mathfrak{J}\to
N\otimes_A A/\mathfrak{J}$~;
la conclusion résulte de ce que $\theta$ et $u$ (resp. $\theta_0$ et
$u_0$) sont simultanément surjectifs
(\hyperref[I.1.3.9-fr]{1.3.9} et
\hyperref[I.10.10.2-fr]{10.10.2}) et de ce que $u$ et $u_0$ sont
simultanément surjectifs (\hyperref[0.7.1.14-fr]{0, 7.1.14}).

\begin{env}[10.11.6]
\label{I.10.11.6-fr}
Le th. (\hyperref[I.10.11.3-fr]{10.11.3}) montre qu'on peut considérer
tout $\mathscr{O}_{\mathfrak{X}}$-Module cohérent $\mathscr{F}$ comme
un $\mathscr{O}_{\mathfrak{X}}$-Module \emph{topologique}, en le
considérant comme limite projective des faisceaux de groupes
\emph{pseudo-discrets} $\mathscr{F}_n$
(\hyperref[0.3.8.1-fr]{0, 3.8.1}). Il résulte alors de
(\hyperref[I.10.11.4-fr]{10.11.4}) que tout homomorphisme
$u:\mathscr{F}\to\mathscr{G}$ de
$\mathscr{O}_{\mathfrak{X}}$-Modules cohérents est automatiquement
\emph{continu}

